% from ega4-16-fr.tex
\oldpage[IV-4]{31}
changement de base $g:S'\to S$, et tout point
$x'\in X'=X\times_S S'$ au-dessus de $x$, on a un isomorphisme canonique
\[
  T_{X'/S'}(x')\xrightarrow{\sim}
  T_{X/S}(x)\otimes_{k(x)}k(x').
  \tag{16.5.13.2}
  \label{IV.16.5.13.2}
\]

Si $x$ est \emph{rationnel sur $k(s)$}, où $s=f(x)$ (de sorte que
$k(s)\to k(x)$ est un isomorphisme), il résulte de (16.4.13) que l'on a un
isomorphisme canonique
\[
  T_{X/S}(x)=T_{X_s/k(s)}(x)
  =\operatorname{Hom}_{k(s)}
  (\mathfrak{m}'_x/\mathfrak{m}'_x{}^2,k(s))
  \tag{16.5.13.3}
  \label{IV.16.5.13.3}
\]
où $\mathfrak{m}'_x$ est l'idéal maximal de
$\mathcal{O}_{X_s,x}=\mathcal{O}_{X,x}/
\mathfrak{m}_s\mathcal{O}_{X,x}$. Dans le cas où $S$ est le spectre d'un corps
$k$, on retrouve ainsi la définition de Zariski de l'espace tangent en un
point $x\in X$ \emph{rationnel sur $k$}, comme dual de
$\mathfrak{m}_x/\mathfrak{m}_x^2$.

Soit $Y$ un second $S$-préschéma et soit $g:Y\to X$ un $S$-morphisme; on a
alors un homomorphisme canonique de $\mathcal{O}_Y$-Modules (16.4.19)
\[
  g_{Y/X/S}:g^*(\Omega_{X/S}^1)\longrightarrow\Omega_{Y/S}^1.
  \tag{16.5.13.4}
  \label{IV.16.5.13.4}
\]
Notons maintenant que si $y\in Y$ et $x=g(y)$, on a
\[
  g^*(\Omega_{X/S}^1)\otimes_{\mathcal{O}_Y}k(y)
  =(\Omega_{X/S}^1\otimes_{\mathcal{O}_X}k(x))
  \otimes_{k(x)}k(y)
\]
et par suite, si $\Omega_{X/S}^1$ est un $\mathcal{O}_X$-Module
\emph{de type fini}, on peut identifier
\[
  \operatorname{Hom}_{k(y)}
  (g^*(\Omega_{X/S}^1)\otimes_{\mathcal{O}_Y}k(y),k(y))
\]
à $T_{X/S}(x)\otimes_{k(x)}k(y)$. On déduit donc alors de l'homomorphisme
(16.5.13.4) un homomorphisme de \emph{$k(y)$-espaces vectoriels}
\[
  T_y(g):T_{Y/S}(y)\longrightarrow
  T_{X/S}(x)\otimes_{k(x)}k(y)
  \tag{16.5.13.5}
  \label{IV.16.5.13.5}
\]
dit \emph{application linéaire tangente à $g$ au point $y$}. Lorsque $y$ est
\emph{rationnel sur $k(s)$}, on peut identifier $k(s)$, $k(y)$ et $k(x)$ et
$T_y(g)$ est alors un homomorphisme de $k(s)$-espaces vectoriels
$T_{Y/S}(y)\to T_{X/S}(x)$; on notera d'ailleurs que dans ce cas,
$g^*(\Omega_{X/S}^1)\otimes_{\mathcal{O}_Y}k(y)$ s'identifie à
$\Omega_{X/S}^1\otimes_{\mathcal{O}_X}k(x)$, et l'homomorphisme précédent est
donc défini sans aucune condition de finitude sur $\Omega_{X/S}^1$ et n'est
autre que l'homomorphisme $T_{Y/S}(g)$ (16.5.12) restreint à la fibre au point
$y$ de $T_{Y/S}$.
\end{env}

\begin{env}[16.5.14]
\label{IV.16.5.14}
L'interprétation des dérivations d'une $A$-algèbre $B$ dans un $B$-module
$L$, donnée dans (0, 20.1.1), se traduit dans le langage des préschémas de la
façon suivante.

Considérons deux morphismes de préschémas $f:X\to S$, $g:Y\to S$, et un
sous-préschéma fermé $Y_0$ de $Y$ défini par un Idéal de carré nul
$\mathscr{J}$ de $\mathcal{O}_Y$ (de sorte que $Y$ et $Y_0$ ont
\emph{même espace topologique sous-jacent}). Supposons donné un $S$-morphisme
$u_0:Y_0\to X$, de sorte qu'on a un diagramme commutatif
\[
  \xymatrix{
    X\ar[d]_-f & Y_0\ar[l]_-{u_0}\ar[d]^j\\
    S & Y\ar@{-->}[ul]_-u\ar[l]^-g
  }
  \tag{16.5.14.1}
  \label{IV.16.5.14.1}
\]
